\chapter[Proposed Evidence-Based Trust Framework]{Proposed Evidence-Based\\ Trust Framework}
\label{ch:master-framework}

\section{Introduction}

The previous chapter showed that Android third-party library detection has
advanced from package-name heuristics toward structural, graph-based,
optimization-aware, and learning-assisted methods. These techniques are
valuable, but they do not by themselves answer the trust question required by
an \SBOM: when a compiled \APK yields a component name, a candidate version,
or a vulnerability warning, what exactly has been proven, what remains
uncertain, and how safely may the result be used?

This chapter proposes an evidence-based trust framework for reconstructed
\SBOM claims. The framework is not a replacement for existing detection
methods such as dependency-graph matching, opcode similarity, learned Smali
features, or source-known validation \cite{zhan2021systematic,zhan2021atvhunter,pan2025libattention,zhou2025sad}.
Instead, it defines a layer above those methods. That layer converts detector
outputs into explicit claims with evidence, provenance, confidence boundaries,
ambiguity, and permitted downstream interpretation.

The central idea is simple: every recovered component entry should be treated
as a claim, not as an unqualified fact. A claim may be exact, bounded,
identity-only, report-only, or out-of-scope. It may be supported by direct
bytecode evidence, by independent structural evidence, by source validation,
or only by contextual hints. It may also be limited by missing references,
partial inclusion, R8 optimization, dependency shading, native code, or hidden
runtime loading. The framework makes these distinctions visible so that an
\SBOM user can understand not only what was emitted, but why it was emitted
and what it is allowed to mean.

\section{Framework Design Principles}

The framework is built around ten design principles. They are intentionally
conservative because reconstructed \SBOM{}s are often used for security
decisions, compliance checks, and vulnerability triage.

\begin{enumerate}[label=\textbf{P\arabic*.},leftmargin=2.2cm]
  \item \textbf{Claim explicitness.}
  Every emitted component, version, relationship, or vulnerability finding is
  represented as a claim with an identifier, a claim type, supporting evidence,
  and a current state. The pipeline does not silently transform a similarity
  score into a fact.

  \item \textbf{Identity-version separation.}
  Component identity and exact version are separate claims. Evidence that
  supports the presence of \texttt{okhttp} or \texttt{androidx.core} does not
  automatically support a specific version. This separation is essential for
  \CVE interpretation because most advisories are version-dependent.

  \item \textbf{Evidence provenance.}
  Each signal records where it came from: \APK-side bytecode, manifest
  metadata, resource evidence, reference database profiles, source-assisted
  validation, external package metadata, vulnerability databases, or
  report-only model output.

  \item \textbf{Reference coverage awareness.}
  A negative or weak result is interpreted in relation to the available
  reference database. Missing coordinates, missing versions, failed feature
  extraction, and incomplete profile families are treated as coverage risks,
  not as proof that a component is absent.

  \item \textbf{Transformation-bounded confidence.}
  Confidence is capped when the analyzed artifact is affected by transformations
  that the evidence cannot reliably handle, including package renaming, code
  shrinking, partial inclusion, repackaging, multidex splitting, class merging,
  and method inlining \cite{androidShrinkCode,xie2024libhunter}.

  \item \textbf{Conflict preservation.}
  Competing versions, adjacent library families, or disagreeing evidence
  sources remain visible. The framework prefers a bounded ambiguity set or an
  abstention over a false sense of exactness.

  \item \textbf{Validation-lane separation.}
  Controlled fixtures, clean source-validated real applications, partial
  source-validated applications, source-known upper bounds, diagnostic
  experiments, and limitation cases are separate lanes. Scores from one lane
  must not be used as if they prove another lane.

  \item \textbf{Authority monotonicity.}
  Downstream authority cannot exceed upstream evidence. A vulnerability claim
  cannot be stronger than the component and version claim on which it depends;
  a standard \CycloneDX or \SPDX entry cannot hide unresolved ambiguity behind
  a syntactically valid package URL \cite{spdxSpecification,cyclonedxSpecification,packageUrlSpecification}.

  \item \textbf{Report-only containment.}
  Diagnostic or machine-learning assistance may guide investigation, ranking,
  and future work, but it is not allowed to modify authoritative component,
  version, or vulnerability outputs until it passes the same validation and
  regression gates as the deterministic evidence.

  \item \textbf{Auditable abstention.}
  Not emitting a component, emitting identity without version, or emitting a
  version band is a legitimate outcome when evidence is insufficient. The
  framework records why the pipeline abstained so that uncertainty becomes a
  visible result rather than a silent failure.
\end{enumerate}

\section{Recovered Claim Model}

\subsection{Claim Definition}

A recovered claim is a structured statement about a component-related fact in
an analyzed \APK. The claim may concern component identity, exact version,
version range, artifact family, dependency relationship, file evidence, or
vulnerability exposure. Formally, a claim can be described as:

\[
  C = \langle s, t, a, e, p, r, q, u, v, \alpha \rangle
\]

where \(s\) is the analyzed subject, \(t\) is the claim type, \(a\) is the
asserted value, \(e\) is the supporting evidence set, \(p\) is provenance,
\(r\) is reference and transformation risk, \(q\) is the evidence quality
profile, \(u\) is uncertainty or ambiguity, \(v\) is validation context, and
\(\alpha\) is the permitted downstream authority. This notation is not meant
to impose a specific storage format. It is a discipline for preventing weak
evidence from being consumed as a stronger result.

\begin{table}[H]
\centering
\small
\caption{Core fields of a reconstructed \SBOM claim.}
\label{tab:master-claim-schema}
\begin{tabularx}{\textwidth}{>{\raggedright\arraybackslash}p{0.20\textwidth}>{\raggedright\arraybackslash}p{0.28\textwidth}X}
\toprule
\textbf{Field} & \textbf{Question answered} & \textbf{Examples}\\
\midrule
Subject & What artifact is being analyzed? & APK identifier, package name,
hash, build variant, DEX or native file scope.\\
Claim type & What kind of statement is this? & Component identity, exact
version, bounded version set, relationship, file component, vulnerability
status.\\
Asserted value & What is being claimed? & Maven coordinate, \PURL, version,
candidate family, dependency edge, native file hash.\\
Evidence set & Why is the claim supported? & Class dependency graph,
opcode/literal features, package metadata, reference profile match, source
validation, advisory match.\\
Provenance & Where did the evidence come from? & APK-side extraction,
reference database, source repository validation, external vulnerability
source, report-only model.\\
Coverage context & What references were available? & Present or missing
coordinate versions, artifact-family coverage, extraction failures, source
resolution status.\\
Transformation risk & What may have damaged the signal? & R8 shrinking,
method inlining, package flattening, partial inclusion, multidex, hidden DEX,
native payloads.\\
Quality profile & How reliable is the evidence? & Directness, uniqueness,
independence, reproducibility, version specificity.\\
Uncertainty & What is unresolved? & Ambiguous adjacent versions, family-level
match only, missing-source component, unsupported native package version.\\
Validation context & Where has the claim been checked? & Controlled fixture,
clean real APK, partial real APK, source-known upper bound, diagnostic lane.\\
Permitted authority & How may it be used? & Confirmed SBOM component,
potential component, report-only hint, confirmed vulnerability, potential
vulnerability, out-of-scope boundary.\\
\bottomrule
\end{tabularx}
\end{table}

\subsection{Claim Granularity}

The framework distinguishes four levels of component granularity:

\begin{enumerate}
  \item \textbf{Identity claim:} a component family or coordinate is present,
  but the exact version is not established.
  \item \textbf{Version claim:} a specific version or bounded version set is
  supported by version-discriminating evidence.
  \item \textbf{Relationship claim:} the component is direct, transitive,
  bundled, copied, shaded, partially included, or uncertain.
  \item \textbf{Security claim:} a vulnerability advisory applies, may apply,
  or does not apply under the current evidence.
\end{enumerate}

This granularity is important because \SBOM formats can represent package
entries even when the recovery process is uncertain. The framework therefore
separates syntactic representation from epistemic authority. A package URL can
identify a candidate component, but the confidence class explains whether the
\PURL is exact, bounded, or only investigatory.

\section{Evidence Hierarchy}

The evidence hierarchy ranks evidence by what it can safely support. It is not
a universal probability scale. Instead, it is a policy for deciding which
interpretations are permitted when the available evidence is incomplete or
conflicted.

\begingroup
\small
\setlength{\LTpre}{0.5em}
\setlength{\LTpost}{0.5em}
\begin{longtable}{>{\raggedright\arraybackslash}p{0.10\textwidth}>{\raggedright\arraybackslash}p{0.22\textwidth}>{\raggedright\arraybackslash}p{0.30\textwidth}>{\raggedright\arraybackslash}p{0.24\textwidth}}
\caption{Evidence hierarchy for reconstructed component claims.}
\label{tab:master-evidence-hierarchy}\\
\toprule
\textbf{Level} & \textbf{Meaning} & \textbf{Typical evidence} & \textbf{Permitted interpretation}\\
\midrule
\endfirsthead
\toprule
\textbf{Level} & \textbf{Meaning} & \textbf{Typical evidence} & \textbf{Permitted interpretation}\\
\midrule
\endhead
E0 & No usable component evidence & No matching code, missing extraction,
unsupported payload type, or unavailable reference family & No package-level
claim. The result may record a boundary or extraction failure.\\
E1 & Weak contextual evidence & Package-name fragment, manifest hint,
resource string, file name, or nearby ecosystem cue & Investigation hint only;
not enough for an \SBOM component claim.\\
E2 & Identity-supporting evidence & Class-set overlap, namespace pattern,
partial structural match, or learned presence signal & Candidate component
identity; exact version and vulnerability authority remain unavailable.\\
E3 & Independent multi-source identity evidence & Agreement between structural
graph evidence, opcode/literal evidence, reference family metadata, or
independent detector families & Strong identity claim, but exact version still
requires version-discriminating evidence.\\
E4 & Version-discriminating APK evidence & Version-specific graph, opcode,
literal, API, or stateful/stateless class features that separate nearby
versions & Exact version or bounded version set, subject to conflict,
coverage, and transformation caps.\\
E5 & Externally validated exact evidence & Source-resolved dependency tree,
controlled fixture ground truth, reproducible build evidence, or manually
verified exact-version evidence & Confirmed exact claim within the validation
context; still not proof of global detector completeness.\\
\bottomrule
\end{longtable}
\endgroup

Two aspects of the hierarchy are deliberately strict. First, evidence levels
are claim-specific. A signal may be E3 for identity but only E1 for exact
version. Second, external validation does not erase the conditions under which
it was obtained. A controlled fixture can validate pipeline mechanics, while a
source-validated real application can validate emitted versions for that
application. Neither automatically proves that all real-world dependencies
are recoverable.

\section{Evidence Quality Dimensions}

Evidence level describes what a signal may support. Evidence quality describes
how fragile that support is. The framework evaluates each evidence item along
the dimensions in \cref{tab:master-quality-dimensions}.

\begin{table}[H]
\centering
\small
\caption{Evidence quality dimensions used to bound claim confidence.}
\label{tab:master-quality-dimensions}
\begin{tabularx}{\textwidth}{>{\raggedright\arraybackslash}p{0.23\textwidth}>{\raggedright\arraybackslash}p{0.34\textwidth}X}
\toprule
\textbf{Dimension} & \textbf{Trust question} & \textbf{Examples of risk}\\
\midrule
Directness & Does the signal observe the component itself or only context
around it? & Manifest strings and package names are indirect; class and method
structure are more direct.\\
Uniqueness & Does the signal distinguish one component or version from close
neighbors? & Shared utility classes, shaded code, and adjacent AndroidX
versions reduce uniqueness.\\
Independence & Do multiple signals fail together or independently? & Two
features derived from the same package namespace are weaker than namespace
plus graph plus literal evidence.\\
Version specificity & Can the evidence separate exact versions? & Identity
features may confirm a coordinate while still leaving a version band.\\
Transformation resilience & Does the signal survive obfuscation and
optimization? & R8 shrinking, inlining, class merging, and package flattening
can remove or distort features.\\
Reference coverage & Were the relevant reference versions available and
successfully extracted? & Missing historical versions or failed profiles cap
confidence even when the best available candidate scores highly.\\
Provenance reliability & Is the evidence reproducible and auditable? &
Manually curated notes, source-resolution logs, and generated feature
artifacts carry different audit strength.\\
Conflict level & Do plausible alternatives remain? & Near-ties, artifact
family ambiguity, and disagreement between source and APK-side evidence require
explicit conflict handling.\\
\bottomrule
\end{tabularx}
\end{table}

These dimensions allow a detector to be useful without being overconfident.
For example, a learned Smali classifier may provide transformation-resilient
identity evidence, but if it lacks exact-version specificity and independent
validation, the framework keeps it report-only or identity-only. Conversely,
source validation can strongly confirm emitted versions for one real
application, but if the source tree contains dependencies that do not appear in
the compiled artifact, the result remains an emitted-version confirmation
rate rather than full dependency recall.

\section{Evidence Fusion, Conflict, and Abstention}

\subsection{Fusion Policy}

The framework treats evidence fusion as an auditable decision process rather
than a single opaque score. A practical implementation may still use numeric
similarity scores, classifier outputs, and thresholds, but the final claim is
assigned by policy. Formal evidence theories provide richer mathematical
fusion rules \cite{shafer1976evidence}; the present framework does not claim to
implement Dempster--Shafer inference and instead uses an explicit,
provenance-preserving authority policy:

\[
  \text{trust}(C) = f(S,D,I,V) - g(K,R,T)
\]

This expression is conceptual. It should not be read as a calibrated
probability unless a separate calibration study has been performed. Its role
is to make the monotonicity rule explicit: more independent, direct, and
version-specific evidence can increase trust, while unresolved conflict,
missing references, and transformation risk can only cap or reduce it.
This caution is necessary because even high-performing learned scores can be
poorly calibrated \cite{guo2017calibration}.
Here, \(S\), \(D\), \(I\), and \(V\) denote support, directness, independence,
and version specificity; \(K\), \(R\), and \(T\) denote conflict, reference
coverage risk, and transformation risk.

\begin{figure}[H]
\centering
\begin{tikzpicture}[
  node distance=8mm and 8mm,
  stage/.style={
    rectangle,
    rounded corners,
    draw=schoolblue,
    thick,
    fill=blue!5,
    align=center,
    minimum width=30mm,
    minimum height=10mm
  },
  risk/.style={
    rectangle,
    rounded corners,
    draw=warningorange,
    thick,
    fill=orange!8,
    align=center,
    minimum width=30mm,
    minimum height=10mm
  },
  claimout/.style={
    rectangle,
    rounded corners,
    draw=evidencegreen,
    thick,
    fill=green!7,
    align=center,
    minimum width=30mm,
    minimum height=10mm
  },
  arrow/.style={-{Latex[length=2.2mm]},thick,draw=schoolblue}
]
\node[stage] (extract) {APK Evidence\\Extraction};
\node[stage,right=of extract] (reference) {Reference and\\Source Context};
\node[stage,right=of reference] (quality) {Quality and\\Provenance};
\node[risk,below=of quality] (risk) {Coverage, Conflict,\\Transformation Caps};
\node[claimout,left=of risk] (claim) {Claim Class and\\Ambiguity};
\node[claimout,left=of claim] (authority) {SBOM and\\Security Authority};

\draw[arrow] (extract) -- (reference);
\draw[arrow] (reference) -- (quality);
\draw[arrow] (quality) -- (risk);
\draw[arrow] (risk) -- (claim);
\draw[arrow] (claim) -- (authority);
\draw[arrow] (extract.south) |- (claim.north);
\end{tikzpicture}
\caption{Conceptual flow of the proposed evidence-based trust framework.}
\label{fig:master-framework-overview}
\end{figure}

\subsection{Conflict Handling}

Conflict is not an exceptional case in APK-side \SBOM recovery; it is normal.
Adjacent versions may share most bytecode, libraries may be partially bundled,
old source trees may resolve different dependency versions from current
reference repositories, and optimized builds may remove the features that
would distinguish two candidates. The framework therefore defines the
following conflict actions:

\begin{enumerate}
  \item \textbf{Exact promotion:} promote a single exact version only when the
  selected candidate has version-discriminating evidence and plausible
  alternatives are ruled out by independent signals or validation.
  \item \textbf{Bounded ambiguity:} emit a version set or band when the
  evidence supports a component family but cannot safely separate nearby
  versions.
  \item \textbf{Identity-only output:} emit the component identity without a
  version when version evidence is weak, missing, or transformation-damaged.
  \item \textbf{Report-only output:} preserve diagnostic candidates for human
  review or future experiments without changing authoritative \SBOM fields.
  \item \textbf{Abstention:} record no component claim when the evidence does
  not meet the minimum identity threshold or when the payload is outside the
  method's scope.
\end{enumerate}

\subsection{Confidence Caps}

Confidence capping is the key mechanism that prevents overclaiming. A claim
cannot exceed the weakest relevant boundary. For example, if the best
available reference database lacks the source version, an exact-version match
against a newer nearby reference may remain a candidate rather than a
confirmed result. If a detector has only report-only validation, its result
cannot modify authoritative \CycloneDX components. If a native shared object is
observed but no native package-version resolver exists, the framework may
emit a file component with hash evidence but not a Maven or native package
version claim.

\section{Claim Classes}

\Cref{tab:master-claim-classes} defines the claim classes used by the
framework. They are ordered by permitted authority, not by how useful they are.
An identity-only or bounded claim can be more honest, and therefore more
useful, than a false exact claim.

\begin{table}[H]
\centering
\small
\caption{Claim classes and permitted downstream use.}
\label{tab:master-claim-classes}
\begin{tabularx}{\textwidth}{>{\raggedright\arraybackslash}p{0.22\textwidth}>{\raggedright\arraybackslash}p{0.31\textwidth}X}
\toprule
\textbf{Claim class} & \textbf{Required support} & \textbf{Permitted use}\\
\midrule
\claimclass{Confirmed exact} & E5 validation or E4 evidence with strong
independent corroboration and no unresolved nearby conflict & Authoritative
\SBOM component and exact-version-dependent security interpretation.\\
\claimclass{Strong exact} & E4 version-discriminating evidence with adequate
reference coverage and manageable transformation risk & Authoritative
component entry with exact version, but validation context must be shown.\\
\claimclass{Likely exact} & E4 evidence exists, but one quality dimension is
weak or a limited alternative remains & Component entry may be emitted with
confidence caveat; vulnerability findings should be conservative.\\
\claimclass{Version-band ambiguous} & Identity is strong, but exact version
cannot be separated from nearby candidates & Emit component with bounded
version set or ambiguity field; exact \CVE classification is not allowed
without additional policy.\\
\claimclass{Identity-only} & E2 or E3 identity evidence without
version-discriminating support & Emit candidate component identity; do not
emit exact-version \CVE conclusions.\\
\claimclass{Report-only candidate} & Diagnostic, experimental, or
machine-learning evidence that has not passed authority gates & Human
investigation, prioritization, or future work only.\\
\claimclass{Unsupported or abstained} & Evidence below identity threshold or
conflict too severe & No package claim; record reason for abstention if useful.\\
\claimclass{Boundary or out-of-scope} & Payload is observed, but the method
does not support its package or version semantics & File/hash component,
limitation note, or risk flag rather than package-version recovery.\\
\bottomrule
\end{tabularx}
\end{table}

\section{Validation and Evidence Lanes}

Evaluation results are meaningful only when their lane is explicit. The
framework separates the lanes in \cref{tab:master-validation-lanes} so that
strong evidence in one context is not accidentally generalized to another.

\begingroup
\small
\setlength{\LTpre}{0.5em}
\setlength{\LTpost}{0.5em}
\begin{longtable}{>{\raggedright\arraybackslash}p{0.21\textwidth}>{\raggedright\arraybackslash}p{0.32\textwidth}>{\raggedright\arraybackslash}p{0.31\textwidth}}
\caption{Validation and evidence lanes.}
\label{tab:master-validation-lanes}\\
\toprule
\textbf{Lane} & \textbf{What it can show} & \textbf{What it must not imply}\\
\midrule
\endfirsthead
\toprule
\textbf{Lane} & \textbf{What it can show} & \textbf{What it must not imply}\\
\midrule
\endhead
Controlled ground truth & Pipeline mechanics, exact expected versions,
known transformations, and regression behavior under designed fixtures & Not
a measure of real-world dependency completeness.\\
Clean source-validated real APKs & Exact emitted Maven versions can be checked
against source-resolved dependency evidence for real applications & Not a
global real-world score if the clean set is small or sparse.\\
Partial source-validated real APKs & Harder real applications can be analyzed
with exact matches, mismatches, missing-source rows, and ambiguity separated &
Partial evidence is not a total failure, but it is also not full recall.\\
Source-known upper bound & Source trees can indicate dependencies that may or
may not be compiled into the APK & Source dependency presence alone does not
prove APK-side presence.\\
Report-only diagnostic lane & Experimental signals can be inspected without
changing authoritative results & Diagnostic improvement is not permission to
alter \SBOM or \CVE outputs.\\
Experiment-only optimization lane & Narrow fixes can be tested with controls,
near-tie analysis, and regression checks & A scoped preflight is not a shared
default until controls and evidence gates pass.\\
Limitation or boundary lane & Native payloads, hidden code, unsupported
transformations, and extraction failures can be reported honestly & Boundary
cases must not be counted as Maven-version failures or successes.\\
\bottomrule
\end{longtable}
\endgroup

This lane structure also changes how metrics should be reported. A controlled
fixture success rate, a clean real-\APK emitted-version confirmation rate, and
a partial real-\APK confirmation rate answer different questions. Combining
them into one number can be useful only after the meaning and denominator are
made explicit. In particular, an emitted Maven version confirmation rate
should not be described as full dependency recall unless the evaluation has
shown that all relevant source dependencies are present in the compiled
artifact and within scope.

\section{Vulnerability-Claim Policy}

\CVE and advisory mapping is downstream of component recovery. The framework
therefore applies an authority rule: vulnerability status cannot be stronger
than the recovered component and version evidence. This is especially
important when advisory ranges are narrow, when adjacent versions have
different security status, or when the result is identity-only. The policy is
compatible with VEX-style separation between component inventory and the
status or applicability of a vulnerability claim \cite{cyclonedxVex}.

\begin{table}[H]
\centering
\small
\caption{Vulnerability interpretation policy.}
\label{tab:master-vulnerability-policy}
\begin{tabularx}{\textwidth}{>{\raggedright\arraybackslash}p{0.25\textwidth}>{\raggedright\arraybackslash}p{0.35\textwidth}X}
\toprule
\textbf{Component evidence} & \textbf{Allowed vulnerability status} & \textbf{Rationale}\\
\midrule
Confirmed exact vulnerable version & Confirmed affected & The version claim
has sufficient authority to evaluate version-specific advisory ranges.\\
Confirmed exact non-affected version & Confirmed not affected for that
advisory & The result is advisory-specific; absence of one advisory does not
prove the component is risk-free.\\
Likely exact version & Potential affected or likely affected, depending on
policy & The version may be used for triage, but the confidence caveat remains
visible.\\
Version-band ambiguous & Potential affected if any candidate in the band is
affected & Exact affected/not-affected labels are unsafe unless all candidates
share the same advisory status.\\
Identity-only component & Potential exposure only when advisory identity
matches and version is unknown & Version-specific claims are not permitted.\\
Report-only candidate & Report-only vulnerability hint & The finding may guide
investigation but must not change authoritative security status.\\
Boundary or out-of-scope payload & File/hash risk note or external analysis
required & Package-version advisory mapping is unavailable in the current
method.\\
\bottomrule
\end{tabularx}
\end{table}

This policy supports conservative security reporting. It avoids two common
errors: hiding a possible issue because the exact version is unknown, and
declaring a confirmed vulnerability from evidence that only supports component
identity. The result can still be actionable: a potential finding may trigger
manual review, dependency upgrade investigation, or source validation, while a
confirmed finding can support stronger remediation language.

\section{Recovered Claim Lifecycle}

The framework can be implemented as a claim lifecycle with explicit state
transitions:

\begin{enumerate}
  \item \textbf{Extract evidence} from the APK, including bytecode,
  resources, manifests, native files, and available build metadata.
  \item \textbf{Generate candidates} from reference profiles, structural
  features, learned hints, package metadata, or source-known context.
  \item \textbf{Normalize identifiers} into coordinates, package URLs, file
  hashes, advisory identifiers, and relationship candidates.
  \item \textbf{Attach provenance} for every evidence item and candidate
  source.
  \item \textbf{Assess quality and coverage} by checking directness,
  independence, version specificity, reference availability, and transformation
  exposure.
  \item \textbf{Fuse evidence} while preserving conflicts, near-ties, and
  missing-reference explanations.
  \item \textbf{Assign claim class} with confidence, ambiguity, validation
  lane, and permitted downstream authority.
  \item \textbf{Emit or abstain} in the \SBOM, vulnerability report, and
  diagnostic artifacts according to the assigned authority.
\end{enumerate}

The lifecycle is intentionally auditable. A reviewer should be able to trace a
\CycloneDX component entry or vulnerability row back to the features, reference
profiles, source validation, and policy decisions that produced it.

\section{PFE System as a Case Study}

The PFE implementation associated with this thesis provides a concrete case
study for the framework. The goal of this section is not to present final
implementation evaluation as a second master experiment. Instead, the mapping
shows how the proposed trust model organizes evidence from a real
APK-to-\SBOM pipeline, and the following section evaluates whether the model
preserves the project's established evidence lanes and authority boundaries.

\begingroup
\small
\setlength{\LTpre}{0.5em}
\setlength{\LTpost}{0.5em}
\begin{longtable}{>{\raggedright\arraybackslash}p{0.25\textwidth}>{\raggedright\arraybackslash}p{0.29\textwidth}>{\raggedright\arraybackslash}p{0.34\textwidth}}
\caption{Mapping PFE observations to framework concepts.}
\label{tab:master-case-study-map}\\
\toprule
\textbf{Case-study observation} & \textbf{Framework concept} & \textbf{Safe interpretation}\\
\midrule
\endfirsthead
\toprule
\textbf{Case-study observation} & \textbf{Framework concept} & \textbf{Safe interpretation}\\
\midrule
\endhead
Controlled transformed fixtures with known expected Maven versions &
Controlled ground-truth lane and E5 validation & Demonstrates pipeline
mechanics and transformation-specific regressions under designed conditions;
does not prove complete real-world recovery.\\
Clean source-validated real APKs with exact emitted Maven-version checks &
Clean real-\APK validation lane & Confirms emitted versions for those
applications and denominators; should remain separate from partial and sparse
cases.\\
Partial source-validated real APKs with exact matches, version mismatches, and
missing-source rows separated & Partial real-\APK lane, conflict preservation,
and emitted-version confirmation rate & Shows useful recovery in difficult
applications without presenting unresolved rows as either success or total
failure.\\
Reference coverage and matcher recall measured as separate gaps & Reference
availability, candidate-generation coverage, and denominator separation &
Shows whether failure arises because the true release is absent from the
reference base or because available evidence was not selected; neither gap
should be hidden inside one accuracy value.\\
Source-known constrained reruns with exact source versions available &
Source-known upper-bound lane & Demonstrates recognition under a favorable
candidate space; does not measure unconstrained APK-only accuracy or prove
that every source dependency survives in the binary.\\
Version-band ambiguity and adjacent-version near-ties & Bounded ambiguity,
confidence caps, and experiment-only optimization lane & Prevents a weak
near-tie from becoming a confirmed exact version or a confirmed \CVE status.\\
Broad scoring and reference-database changes that improved targets but
regressed controls & Negative evidence and promotion governance & Preserves
failed experiments as evidence that a signal is unsafe for shared authority,
even when it is useful diagnostically.\\
The scoped AndroidX Fragment 1.6.2 preflight selected the intended release
without sampled control regressions, but lacked version-unique APK evidence &
Experiment-only near-tie decision & Keeps the result outside the shared
reference database because clean sampled controls do not convert
feature-equivalent candidates into exact evidence.\\
Report-only R8 and AI assistance & Report-only containment & Allows learned
or diagnostic signals to support triage while keeping authoritative \SBOM and
vulnerability outputs deterministic until validation gates pass.\\
Native shared objects and hidden-code cases & Boundary or out-of-scope claim
class & Supports file/hash evidence, limitation notes, or risk flags without
inventing Maven or native package-version claims.\\
\bottomrule
\end{longtable}
\endgroup

This mapping illustrates the main contribution of the framework: it turns
mixed evidence into bounded statements. A controlled result can be strong
without being overgeneralized. A partial real-\APK result can be valuable
without being hidden as a failure. A source-known dependency can inform an
upper bound without proving APK-side presence. A report-only classifier can
guide investigation without becoming \SBOM authority. These distinctions make
the pipeline more cautious, but they also make it more usable because users
can decide how to act on each claim class.

\section{Framework Case-Study Evaluation}

The framework was evaluated against the audited PFE evidence package at two
levels: representational coverage and authority preservation. Representational
coverage asks whether every major project outcome can be expressed without
collapsing its validation context. Authority preservation asks whether weak,
conflicting, source-constrained, or learned evidence remains prevented from
silently changing authoritative component, version, or vulnerability claims.
The underlying ledgers are retained in the project artifacts, including
\path{docs/evaluation/THESIS_EVALUATION_TABLES.md},
\path{docs/thesis/FINAL_V2_RESULTS_SECTION.md}, and the experiment-specific
diagnostic reports.

\begingroup
\small
\setlength{\LTpre}{0.5em}
\setlength{\LTpost}{0.5em}
\begin{longtable}{>{\raggedright\arraybackslash}p{0.21\textwidth}>{\raggedright\arraybackslash}p{0.36\textwidth}>{\raggedright\arraybackslash}p{0.31\textwidth}}
\caption{Case-study evaluation of representational coverage and authority preservation.}
\label{tab:master-framework-evaluation}\\
\toprule
\textbf{Criterion} & \textbf{Observed PFE evidence} & \textbf{Framework evaluation}\\
\midrule
\endfirsthead
\toprule
\textbf{Criterion} & \textbf{Observed PFE evidence} & \textbf{Framework evaluation}\\
\midrule
\endhead
Evidence-lane separation &
The controlled lane confirmed 84/84 expected versions. The clean real lane
confirmed 47/47 emitted versions, while the partial real lane confirmed
880/1790. Four explicit boundary cases were excluded from accuracy metrics. &
Passed for the case study: strong, partial, and boundary evidence remained
separate rather than being pooled into one success claim.\\

Denominator separation &
Across 102 source-validated real APKs, 927/1837 emitted Maven versions were
confirmed. Against the wider 12,082-row source dependency denominator, the
same 927 exact emitted rows represent 7.7\% end-to-end exact source recall. &
Passed: emitted-row correctness and full source-dependency recall are stated
as different metrics with different authority.\\

Upper-bound containment &
In 106 source-known constrained cases, 6772/6772 emitted rows matched the
source versions, while 5238 source rows remained gaps. &
Passed: the result is represented as a favorable source-known recognition
upper bound, not as unconstrained APK-only accuracy.\\

Conflict and negative controls &
The Fragment 1.6.2 preflight selected the intended version in 10/10 target
rows but remained an unresolved near-tie. The coroutines 1.6.4 experiment
fixed 20/20 target rows but regressed three controls. Neither change was
promoted. &
Passed: apparent target improvements did not override ambiguity or negative
control evidence.\\

Report-only containment &
The AI/R8 corpus contains 31,589 rows, with 20,523 training-eligible and
11,066 held-out or report-only rows. Report-only integration left normal
outputs unchanged, and stricter APK-only gates produced zero promotion
candidates. &
Passed: learned assistance can support triage while component selection,
\SBOM emission, and vulnerability authority remain unchanged.\\

Transformation-risk containment &
The Phase 4 multi-scale R8 diagnostic found ten possible real improvements,
but also six controlled regression blockers, 121 real regression-risk rows,
and ten clean-real blockers. Global promotion was rejected. &
Passed: transformation-oriented evidence remained diagnostic when its control
matrix did not support shared authority.\\

Boundary representation &
Four native-heavy or obfuscation cases remained outside accuracy metrics.
Seventeen native files were represented as hash-only evidence; the hybrid
Jitsi case validated 5/7 Maven rows while keeping 13 native files hash-only. &
Passed: unsupported package/version claims were not invented for native,
hidden, or framework-specific payloads.\\
\bottomrule
\end{longtable}
\endgroup

The evaluation therefore supports a bounded conclusion: the framework covers
all major evidence outcomes exercised by the PFE and preserved authority in
the observed negative-promotion, report-only, and boundary cases. It does not
establish calibrated confidence probabilities, completeness across all Android
detectors, or independent analyst agreement. Those questions require external
datasets, tools, and reviewers. The case study nevertheless replaces a purely
proposed evaluation with auditable evidence that the framework can describe
real successes, failures, near-ties, upper bounds, and abstentions without
changing their meaning.

\section{Limitations and Research Agenda}

The proposed framework is intentionally evidence-centric, but it has several
limitations.

\begin{enumerate}
  \item \textbf{It does not solve feature extraction by itself.} The framework
  requires detectors, reference profiles, source validation, and vulnerability
  data. Weak upstream evidence still limits the final claims.

  \item \textbf{It does not provide calibrated probabilities.} Confidence
  classes are policy categories unless supported by a separate calibration
  study over representative datasets.

  \item \textbf{It depends on reference and source availability.} Missing
  Maven versions, unavailable historical artifacts, broken build systems, and
  incomplete source trees can cap evidence even when the APK-side matcher is
  strong.

  \item \textbf{It does not fully solve transitive dependency semantics.}
  Reconstructed bytecode can show that code is present, but distinguishing
  direct, transitive, shaded, copied, and build-time-only dependencies remains
  difficult.

  \item \textbf{It leaves native package-version recovery as future work.}
  Native files can be represented as file components with hashes, but mapping
  them to precise upstream package versions requires a different evidence
  base.

  \item \textbf{It requires disciplined reporting.} The framework is only
  useful if evaluation tables, thesis claims, and generated reports keep the
  lanes and authority boundaries visible.
\end{enumerate}

These limitations define a research agenda. Future work should combine
multi-scale bytecode evidence with calibrated confidence models, richer
reference coverage, reproducible source-validation pipelines, native library
package recovery, and benchmark protocols that measure not only exact matches
but also abstention quality and vulnerability-authority correctness.

\section{Conclusion}

This chapter proposed an evidence-based trust framework for reconstructing
\SBOM claims from compiled Android applications. The framework treats each
component, version, relationship, and vulnerability statement as an auditable
claim with evidence, provenance, quality dimensions, validation context,
ambiguity, and permitted downstream authority. Its main purpose is to prevent
weak or partial evidence from being consumed as stronger security knowledge.

The framework also provides a bridge between the literature reviewed in the
previous chapter and the PFE APK-to-\SBOM implementation. Existing detection
techniques supply useful signals; the framework defines how those signals are
bounded, fused, and reported. This makes trustworthy \SBOM recovery less a
question of always producing one exact answer and more a question of producing
the strongest answer that the evidence can honestly support.
